\chapter{Bond Mathematics and Curve Construction}

\begin{objectives}
Price fixed-income cash flows under market conventions, construct discount curves,
measure rate sensitivity, and diagnose interpolation and bootstrapping risk.
\end{objectives}

\section{Present value and compounding}

For deterministic cash flows $C_i$ at dates $T_i$,
\[
P_t=\sum_i C_iP(t,T_i).
\]
If a flat continuously compounded yield $y$ is used,
$P(y)=\sum_iC_i\e^{-y(T_i-t)}$. Market instruments instead use instrument-specific
compounding, day counts, settlement lags, calendars, and accrued-interest rules.

An annualized rate is incomplete without basis. ACT/360 and ACT/365 change accrual;
30/360 changes date counting; business-day adjustment changes payment date.
Production curve code treats conventions as data.

\section{Bootstrapping}

Bootstrapping solves discount factors sequentially from liquid instruments. A par
coupon bond priced at par satisfies
\[
1=c_n\sum_{i=1}^n\delta_iP(0,T_i)+P(0,T_n).
\]
If earlier discounts are known,
\[
P(0,T_n)=
\frac{1-c_n\sum_{i=1}^{n-1}\delta_iP(0,T_i)}
{1+c_n\delta_n}.
\]

Deposits, futures or forward-rate agreements, overnight swaps, term swaps, and bonds
can occupy different curve segments. Multi-curve construction discounts
collateralized cash flows on an overnight curve while projecting each floating tenor
on its own forward curve.

\section{Interpolation}

Quotes exist at discrete maturities while pricing needs all dates. Interpolating
zero rates, log discounts, forwards, or par rates imposes different shapes.
Linear interpolation of log discount factors gives piecewise constant forwards.
Cubic splines are smooth but can oscillate and create negative forward rates.
Monotone-convex methods aim to preserve economically sensible local behavior.

Curve risk includes quote noise, interpolation choice, instrument selection, and
extrapolation. A perfect repricing of bootstrap instruments does not make the
between-knot curve observed.

\section{Duration and convexity}

For cash-flow weights
$w_i=C_i\e^{-yt_i}/P$, Macaulay duration is
$D_M=\sum_iw_it_i$. With continuous compounding, modified duration equals $D_M$:
\[
-\frac1P\frac{\partial P}{\partial y}=D_M.
\]
Convexity is
$\mathcal C=P^{-1}\partial^2P/\partial y^2=\sum_iw_it_i^2$.

Yield duration assumes a parallel flat-yield move. Key-rate duration shocks one
curve node with an interpolation rule. Partial derivatives to market quotes provide
more realistic curve risk but depend on the bootstrap Jacobian.

\begin{proposition}[Positive convexity of an option-free bond]
For fixed positive cash flows under a parallel continuously compounded yield,
$P''(y)>0$.
\end{proposition}
\begin{proof}
$P''(y)=\sum_i t_i^2C_i\e^{-yt_i}>0$ when at least one positive cash flow occurs
after time zero.
\end{proof}

\section{Immunization}

A liability is locally immunized by matching present value and duration with an
asset portfolio. Convexity improves protection for parallel moves. Multiple
key-rate constraints immunize selected curve deformations:
\[
\min_w w^\transpose\Sigma w
\quad\text{s.t.}\quad
A^\transpose w=b,
\]
where rows of $A$ contain value and key-rate exposures. Rebalancing, transaction
costs, nonparallel moves, spread risk, and optionality limit immunization.

\section{Problems}

\begin{exercise}[Core: bootstrap]
Bootstrap annual discount factors from three par rates and compute one-year forward
rates. Check discount monotonicity.
\end{exercise}
\begin{exercise}[Core: duration]
Derive duration and convexity for a zero-coupon bond and compare the approximation
with the exact price change.
\end{exercise}
\begin{exercise}[Advanced: curve Jacobian]
Use the implicit function theorem to derive sensitivities of bootstrapped discount
factors to input par swap rates.
\end{exercise}


\chapter{Interest-Rate Models and Rate Derivatives}

\begin{objectives}
Model short rates and forward rates under risk-neutral and physical measures, price
caps and swaptions, and understand calibration and measure consistency.
\end{objectives}

\section{Short-rate models}

In a short-rate model,
\[
P(t,T)=\E_t^\Q\left[\exp\left(-\int_t^Tr_s\dd s\right)\right].
\]
Vasicek uses Gaussian mean reversion. Cox--Ingersoll--Ross uses
\[
\dd r_t=\kappa(\theta-r_t)\dd t+\sigma\sqrt{r_t}\dd W_t^\Q,
\]
which is affine and nonnegative under boundary conditions. Hull--White extends
Vasicek with time-dependent drift:
\[
\dd r_t=(\theta(t)-a r_t)\dd t+\sigma\dd W_t^\Q,
\]
allowing exact fit to today's initial curve.

Exact initial fit does not identify future dynamics. Mean reversion and volatility
are calibrated to options or estimated under a measure; the physical drift requires
risk-premium specification.

\section{Heath--Jarrow--Morton}

Instead of a short rate, HJM models the entire instantaneous forward curve:
\[
\dd f(t,T)=\alpha(t,T)\dd t+\sigma(t,T)^\transpose\dd W_t^\Q.
\]
No arbitrage restricts drift:
\[
\alpha(t,T)=
\sigma(t,T)^\transpose\int_t^T\sigma(t,u)\dd u
\]
under the basic money-market measure convention. Thus volatility specification
determines risk-neutral drift.

Finite-dimensional Markov representations require special volatility structures.
Direct curve models make it clear that today's curve is a state, not a scalar rate.

\section{Forward measures and caplets}

The simply compounded forward rate for $[T_1,T_2]$ is
\[
L_t(T_1,T_2)=
\frac{1}{\delta}\left(
\frac{P(t,T_1)}{P(t,T_2)}-1
\right).
\]
Under the $T_2$-forward measure it is a martingale in a standard setting. A caplet
payoff at $T_2$ is
$N\delta(L_{T_1}-K)^+$ and Black's formula is
\[
V_t=NP(t,T_2)\delta
[F\Phi(d_1)-K\Phi(d_2)]
\]
for lognormal forward $F=L_t(T_1,T_2)$.

Negative rates motivate shifted lognormal or normal Bachelier quoting. A normal
volatility has rate units; a lognormal volatility has percentage units. Conversion
depends on forward, strike, and maturity.

\section{Swaptions}

A payer swaption grants the right to pay fixed $K$ and receive floating. Under the
annuity numeraire
\[
A_t=\sum_i\delta_iP(t,T_i),
\]
the forward swap rate
\[
S_t=\frac{P(t,T_0)-P(t,T_n)}{A_t}
\]
is a martingale under its swap measure. Black-style pricing treats it as lognormal:
\[
V_t=A_t[S_t\Phi(d_1)-K\Phi(d_2)].
\]

Smile, cash settlement, day count, collateral, and exercise conventions enter real
contracts. Bermudan swaptions depend on the joint curve evolution and optimal
exercise.

\section{Market models}

The LIBOR market model specifies dynamics of a family of forward rates under their
natural measures. Changing to one common simulation measure introduces drift terms
depending on correlations and other forwards. Calibration matches cap and swaption
volatilities but correlation is weakly identified by vanilla prices and critical for
exotics.

Post-benchmark-reform markets use overnight and term reference rates with fallback
and compounding conventions. Mathematical notation should name the actual rate and
accrual mechanism instead of retaining obsolete labels by habit.

\section{Problems}

\begin{exercise}[Core: CIR moments]
Derive conditional mean of the CIR rate. Research the noncentral chi-square transition
law and explain its exact-simulation benefit.
\end{exercise}
\begin{exercise}[Core: caplet measure]
Derive caplet pricing by choosing the $T_2$ bond as numeraire.
\end{exercise}
\begin{exercise}[Advanced: HJM drift]
Derive the HJM drift restriction by requiring discounted bond prices to be
martingales.
\end{exercise}


\chapter{Credit Risk}

\begin{objectives}
Model default with structural and intensity approaches, bootstrap survival curves,
price credit protection, and separate credit, liquidity, and recovery assumptions.
\end{objectives}

\section{Credit events and loss}

Credit risk includes default, downgrade, spread widening, restructuring, and
counterparty deterioration. If exposure at default is EAD, loss fraction is
$LGD=1-R$, and default indicator is $D$, loss is
\[
L=\mathrm{EAD}\times LGD\times D.
\]
Expected loss is not unexpected loss. Dependence among EAD, recovery, and default is
important in stress.

\section{Structural models}

In Merton's model, firm asset value follows a diffusion and equity is a call on
assets with strike equal to promised debt:
\[
E_0=V_0\Phi(d_1)-D\e^{-rT}\Phi(d_2).
\]
Default at maturity occurs when $V_T<D$. Risky debt equals risk-free debt minus a put
on firm assets.

The model connects leverage and asset volatility to credit spreads, but firm asset
value and volatility are latent, default can occur only at specified barriers, and
short-maturity spreads are often too low without jumps or liquidity.

First-passage models default when asset value crosses a barrier. Barrier design and
capital structure become central.

\section{Reduced-form intensity models}

Let default time $\tau$ have intensity $\lambda_t$ so, informally,
\[
\Prob(t<\tau\leq t+\dd t\mid\tau>t,\F_t)=\lambda_t\dd t.
\]
Conditional survival is
\[
\Prob(\tau>T\mid\F_t)
=\E_t\left[
\exp\left(-\int_t^T\lambda_s\dd s\right)
\right]
\]
under appropriate filtration and measure conditions.

With deterministic intensity and recovery paid at default, a risky zero-coupon value
is the sum of survival payment and expected recovery payment. Under zero recovery,
$P^{\mathrm{risky}}(0,T)=P(0,T)S(0,T)$ when rates and default separate appropriately.

\section{Credit default swaps}

A CDS protection buyer pays spread $s$ until maturity or default and receives
$1-R$ upon covered credit event. The par spread equates premium and protection legs:
\[
s\sum_i\delta_iP(0,T_i)S(0,T_i)
\approx
(1-R)\sum_iP(0,T_i)
[S(0,T_{i-1})-S(0,T_i)],
\]
with production valuation adding accrual on default and exact timing.

Given recovery and discount curve, quoted spreads bootstrap survival probabilities
or piecewise intensities. Recovery and intensity are not separately well identified
by CDS spreads alone.

\section{Credit portfolio dependence}

In a one-factor Gaussian latent-variable model,
\[
Z_i=\sqrt{\rho_i}Y+\sqrt{1-\rho_i}\varepsilon_i,
\]
and obligor $i$ defaults if $Z_i$ is below a threshold matching marginal default
probability. Conditional on $Y$, defaults are independent. The common factor creates
default clustering.

Tail loss is extraordinarily sensitive to asset correlation, exposure concentration,
recovery dependence, and model form. Gaussian dependence can understate joint
extremes. Stress testing should include sector and macro scenarios rather than rely
only on a fitted copula.

\section{Migration matrices}

A rating transition matrix $P$ has rows summing to one. Multi-year transition under
time homogeneity is $P^n$. A continuous-time generator $Q$ has nonnegative
off-diagonals, rows summing to zero, and transition matrix $\exp(Qt)$. Not every
empirical discrete matrix is embeddable in a valid generator.

Historical transition probabilities are physical. Pricing spreads reflects
risk-neutral probabilities, recovery, liquidity, capital, and risk premia.

\section{Counterparty exposure}

Netting sets combine positive and negative trade values under enforceable legal
agreements. Variation margin reduces current exposure; initial margin targets future
exposure over margin period of risk. Rehypothecation, thresholds, disputes, and close-out
delay affect residual exposure.

Potential future exposure, expected shortfall of exposure, CVA, and regulatory
exposure are different metrics. A simulation needs joint market factors, collateral
rules, default dependence, and pathwise netting.

\section{Problems}

\begin{exercise}[Core: Merton]
Show that risky debt equals a risk-free bond minus a put. Derive the risk-neutral
default probability and distinguish it from expected loss.
\end{exercise}
\begin{exercise}[Core: CDS bootstrap]
Bootstrap piecewise constant survival probabilities from two annual CDS spreads
under a simplified payment convention.
\end{exercise}
\begin{exercise}[Advanced: wrong-way risk]
Construct a joint model in which exposure and intensity share a factor. Compare CVA
with an independence approximation.
\end{exercise}


\chapter{Securitization and Structured Credit}

\begin{objectives}
Understand cash-flow waterfalls, tranche loss, prepayment and extension risk, and
why structured-credit models are highly sensitive to dependence and legal detail.
\end{objectives}

\section{Pooling and tranching}

A securitization pools asset cash flows and allocates them through a contractual
waterfall. A tranche with attachment $A$ and detachment $D$ has loss at portfolio
loss fraction $L$:
\[
L_{A,D}(L)=
\min\{(L-A)^+,D-A\}.
\]
Equity absorbs first loss; mezzanine absorbs intermediate loss; senior tranches lose
only after subordination is exhausted.

Expected tranche loss is nonlinear in the full portfolio-loss distribution. Two
portfolios with the same expected total loss can have very different senior-tranche
risk.

\section{Dependence and base correlation}

Increasing default correlation moves mass toward very low and very high portfolio
loss. This can reduce equity expected loss at some points while increasing senior
tail loss. Base correlation fits each equity tranche $[0,D]$ with a separate
correlation, creating an interpolation convention rather than one internally
consistent joint distribution.

\section{Mortgage cash flows}

Mortgage-backed securities contain borrower prepayment options. Falling rates can
accelerate refinancing, shortening duration when price would otherwise rise;
rising rates slow prepayment, extending duration when price falls. This negative
convexity drives dynamic hedging demand.

Prepayment depends on incentive, seasoning, burnout, housing turnover, credit,
seasonality, and operational friction. A model needs loan-level or pool variables
available at the valuation date.

\section{Waterfall simulation}

A pathwise engine applies:

\begin{enumerate}
\item scheduled interest and principal;
\item prepayment and default;
\item recovery lag and severity;
\item fees and servicing;
\item interest and principal priority;
\item triggers, reserve accounts, and diversion;
\item cleanup calls and termination.
\end{enumerate}

The legal prospectus is part of the model specification. A mathematically elegant
loss model cannot compensate for an incorrect waterfall.

\section{Stress and model risk}

Structured products combine long horizons, nonlinear allocation, uncertain behavior,
and sparse tail data. Stress unemployment, home prices, rates, refinancing capacity,
recovery, correlation, and servicer performance. Report scenario cash flows and
breakeven assumptions rather than only a single spread or rating.

\section{Problems}

\begin{exercise}[Core: tranche loss]
Plot tranche loss as a function of portfolio loss for equity, mezzanine, and senior
tranches. Identify delta with respect to total loss away from attachment points.
\end{exercise}
\begin{exercise}[Advanced: large homogeneous pool]
Conditional on a Gaussian common factor, derive the limiting default fraction in a
homogeneous portfolio and integrate to obtain expected tranche loss.
\end{exercise}
